\chapter{Risultati e discussione}
\label{chap:risultati}

% Struttura fissata con il relatore il 27 ottobre 2026. Il capitolo provvisorio
% precedente, basato sulla conferma da E-OBS, è in
% backup/2026-10-27_risultati_provvisori_eobs.tex (non compilato). Le fonti dei
% numeri, quando ci saranno, vanno in chapters/risultati_fonti.tex.
% Regola: ogni numero viene dai file della run finale scelta dalla ricerca dei
% parametri; finché la run non è scelta, le sezioni restano vuote.

Il capitolo riporta e discute i risultati del detector sui campi ERA5, nella configurazione scelta dalla ricerca dei parametri. In assenza di etichette, il comportamento del detector è studiato in quattro passi. Il primo riguarda la scelta della configurazione e la convergenza della ricerca. Il secondo esamina le componenti dello score, del generatore e del discriminatore, e il loro legame nello spazio e nel tempo, per stabilire quale parte dello score usare. Il terzo confronta il detector con metodi più semplici, statistici o basati su modelli poco profondi. L'ultimo esamina casi specifici, prima in generale e poi per tipo di evento.

\section{Ricerca dei parametri}
\label{sec:ris-ricerca}

% Descrizione della ricerca in due fasi (controller bayesiano di W&B, poi
% Optuna ripartito dalle stesse esecuzioni), con la figura dell'assestamento:
% obiettivo di ogni prova e migliore fino a quel momento, e percorso delle
% configurazioni nello spazio dei parametri. Figura da generare fra qualche giorno
% dallo studio Optuna era5-gan (outputs/optuna/era5-gan.db sul server).
% Ricordare: la PCA dei quattro parametri tiene il 38% della varianza sulla prima
% componente, quanto darebbe il caso con 21 esecuzioni; mostrare i parametri
% direttamente (scripts/plot_era5_sweep_pca.py) o concordare con il relatore.
\provv{Sezione da scrivere quando la ricerca con Optuna si sarà assestata.}

\section{Correlazione spaziale degli score}
\label{sec:ris-correlazione-spaziale}

% Score del generatore e del discriminatore, prima e dopo la normalizzazione
% climatologica (G, D, G+clim, D+clim; il combinato non è considerato). A istante
% fissato: mappa dei quattro score. Lungo tutto il test: per ogni istante la
% correlazione spaziale (Pearson e Spearman) fra le sei coppie sulle celle, e il
% grafico di come cambia nel tempo. Script: check_era5_score_correlation.py.
\provv{Sezione da scrivere con la run finale.}

\section{Correlazione temporale degli score}
\label{sec:ris-correlazione-temporale}

% Per ogni cella, la correlazione nel tempo fra le sei coppie di score sul test:
% mappe delle correlazioni. Stesso script della sezione precedente.
\provv{Sezione da scrivere con la run finale.}

\section{Perché il generatore rileva e il discriminatore no}
\label{sec:ris-g-d}

% Spiegare, sulla base delle due sezioni precedenti, perché lo score del
% generatore è informativo e quello del discriminatore no (uscita vicina a 0.5,
% score combinato dominato dalla componente del discriminatore). Contesto dalla
% letteratura: in AnoGAN il solo termine del discriminatore ha AUC 0.74 contro
% 0.88 della ricostruzione (Capitolo 2).
\provv{Sezione da scrivere con la run finale.}

\section{Confronto con metodi semplici e modelli poco profondi}
\label{sec:ris-confronti}

% Confronto a parità di segnalazioni (ognuno il proprio 5% più alto): quota delle
% segnalazioni condivise, indice di Jaccard, segnalazioni esclusive, voto a
% maggioranza (compare_era5_detectors.py). Rilevatori, tutti esistenti e applicati
% così come sono, senza riferimenti esterni:
% - solo il punto: distanza dal normale; sette metodi statistici di TSB-AD (PCA,
%   KMeans, CBLOF, MCD, HBOS, COPOD, Isolation Forest; score_era5_shallow.py);
% - il tempo: variazione in tre ore; Shekhar et al. 2001 con vicinato temporale;
% - lo spazio: Shekhar et al. 2001 con vicinato spaziale; Chen et al. 2008
%   multivariato (score_era5_spatial_outliers.py);
% - spazio e tempo: Shekhar et al. 2001 con vicinato spazio-temporale; STGAN.
% KNN, LOF e OCSVM esclusi per il costo su ~600 milioni di punti. I metodi
% spazio-temporali pubblicati (scan statistic, ST-DBSCAN, ...) non si applicano
% così come sono a campi continui multivariati su griglia.
% Testo tolto dal Capitolo 5 il 28 ottobre 2026 (Metriche e Rilevatori senza modello),
% da riscrivere qui:
% \section{Metriche}
% \label{sec:exp-metriche}
%
% L'accordo fra due rilevatori che segnalano gli insiemi di osservazioni $A$ e $B$ è misurato dall'indice di Jaccard,
% \begin{equation}
% J(A,B)=\frac{|A\cap B|}{|A\cup B|},
% \end{equation}
% che vale 1 quando gli insiemi coincidono e 0 quando sono disgiunti. Quando i rilevatori segnalano la stessa quota di osservazioni è riportata anche la quota delle segnalazioni di $A$ che anche $B$ segnala, $|A\cap B|/|A|$, che a parità di segnalazioni coincide con $|A\cap B|/|B|$. Nessuna delle due è una misura di accuratezza: due rilevatori possono concordare ed entrambi sbagliare.
%
% \section{Rilevatori di confronto senza modello}
% \label{sec:exp-baseline}
%
% Il detector è confrontato con tre rilevatori che non usano un modello appreso. Rappresentano i due modi più semplici di dire che il tempo di una cella è insolito, e servono a stabilire che cosa il detector aggiunge a ciò che si ottiene senza addestramento. Sono calcolati sugli stessi dati standardizzati della Sezione~\ref{sec:ingresso-normalizzazione-era5}.
%
% La \emph{distanza dal normale} è l'errore che si commetterebbe rispondendo, in ogni cella, con il valore normale per quel giorno e quell'ora: misura lo scostamento dal clima, e coglie per esempio il caldo o il freddo prolungati. La \emph{variazione in tre ore} è l'errore che si commetterebbe rispondendo con il valore dell'istante precedente: misura il cambiamento rapido, e coglie per esempio il passaggio di un fronte. Il terzo rilevatore prende, cella per cella, il maggiore dei due. In entrambi gli errori ogni variabile è divisa per il proprio errore medio, così che le quindici variabili contino allo stesso modo, e ai tre score è applicata la stessa forma di normalizzazione climatologica dello score di ricostruzione (Sezione~\ref{sec:normalizzazione-climatologica-score}). Per trattarli come il detector, errori medi e climatologia sono stimati sugli stessi anni di validazione, esclusi i valori segnati come insoliti, e restano fissi per il test; le soglie sono fissate allo stesso modo, al 95-esimo e al 99-esimo percentile degli score di validazione. I confronti sono quindi a parità di criterio.
%
% % Codice: rebalance_era5_stgan_scores.py --scale mean --scale-run-dir (validazione)
% % --normal-masks-dir; climatology --reference validation. Run in outputs/good2v
% % (val_* e test_*).
\provv{Sezione da scrivere con la run finale.}

\section{Eventi per tipo e dimensione}
\label{sec:ris-eventi}

% Eventi dello score scelto nell'ipercubo tempo x latitudine x longitudine,
% divisi per dimensione (1, 2-4, 5-19, >19 valori) e per tipo (la variabile
% d'ingresso più lontana dal normale stagionale sul nucleo). Per ogni score:
% percentile mediano e quota segnalata nel nucleo, nel contorno (le 8 celle
% attorno) e lontano dagli eventi (check_era5_event_zones.py). Nessun
% riferimento esterno.
\provv{Sezione da scrivere con la run finale.}

\section{Casi specifici}
\label{sec:ris-casi}

% Comportamento dello score in tre situazioni: dentro un'anomalia estesa, sul
% contorno segnalato dal solo modello, in una zona non anomala. Prima in generale,
% poi per tipi di evento diversi.
\provv{Sezione da scrivere con la run finale.}

\section{Limiti dei risultati}
\label{sec:ris-limiti}

% Da scrivere alla fine, sui risultati effettivi.
\provv{Sezione da scrivere alla fine.}
